% !TEX root = ../main.tex
% chktex-file 46
\chapter{Zusammenfassung}%
\label{sec:conclusion}

Abschließend wird nun zusammengefasst, inwiefern das in~\ref{sec:intro:goals} beschriebene Ziel einer erweiterbaren, parallelisierbaren Wissensgraph-Konstruktion aus Kommunikationsdaten erreicht wurde.
Hierfür werden in~\ref{sec:conclusion:review} die erreichten Teilziele betrachtet.
Anschließend wird in~\ref{sec:conclusion:todo} ein Ausblick auf Möglichkeiten der Weiterentwicklung des Systems gegeben.

\section{Rückblick}%
\label{sec:conclusion:review}

\begin{figure}[h]
	\centering
	\includegraphics[width=0.65\textwidth]{gfx/conclusion/overview.pdf}
	\caption{Grober Aufbau des Konstruktionsverfahrens}\label{fig:conclusion:overview}
\end{figure}
In Kapitel~\ref{sec:text2kg} wurden Lösungen für die drei wesentlichen Teilprobleme Ontologie, Sprachverarbeitung und Grapherweiterung beschrieben.
Die drei Teillösungen bilden zusammen ein Konstruktionsverfahren, welches dem zu Anfang beschriebenen Aufbau entspricht.

\paragraph{Ontologie}
Die Wissensgraphontologie basiert auf Konzeptgraphen.
Im Gegensatz zu den Fakten-Tripel basierten semantischen Netzen, die lediglich beschreiben, welche Beziehungen zwischen Konzepten existieren, kann mittels Konzeptgraphen auch beschrieben werden, welche Beziehungen nicht, oder nur möglicherweise existieren.
Die konstruierten Wissensgraphen haben also eine deutlich höhere Ausdrucksstärke, als z.~B. NELL oder der Google~Knowledge~Graph.
Diese Ausdrucksstärke ist notwendig, um die Komplexität natürlichsprachlicher Aussagen sinnvoll abbilden zu können.

Damit die Konzeptgraphen eine möglichst präzise Semantik haben, muss die Semantik der verwendeten Relationen beschrieben werden.
Es wurde daher ein möglichst minimales Set von Relationen definiert~\tref{sec:text2kg:ontology:pred}, welches zur Beschreibung einer Vielzahl natürlichsprachlicher Aussagen geeignet ist.

Eine der in~\ref{sec:intro:goals} beschriebenen Anforderungen an das Konstruktionsverfahren ist die Erweiterbarkeit, d.~h.\ die prinzipielle Möglichkeit nicht natürlichsprachliche Eingaben einzufügen.
Konkret bedeutet dies, dass z.~B. ein Bild in denselben Wissensgraph eingefügt werden kann, in den auch Textnachrichten eingefügt werden.
Da die gewählten Konzeptgraph-Relationen nicht speziell für natürliche Sprache ausgelegt sind, ist dies prinzipiell möglich;
es gibt keine Relationen, die bestimmte grammatikalische Konstrukte oder Wortarten abbilden.
Wie die Integration anderer Eingabearten im Detail aussehen könnte, war jedoch nicht Teil der Aufgabenstellung und wurde daher nicht näher untersucht.

\paragraph{Sprachverarbeitung}
Für die Transformation einer Nachricht in einen Konzeptgraphen wird das Stanford~CoreNLP Framework mit einer Menge von Transformationsregeln~\tref{sec:text2kg:nlp:transform} kombiniert, welche die von CoreNLP gefundenen Strukturen auf äquivalente Konzeptgraph-Strukturen abbilden.

\paragraph{Grapherweiterung}
Das Einfügen des aus einer Nachricht extrahierten Konzeptgraphen in einen bestehenden Wissensgraphen erfolgt mittels PSL.\@
Im Rahmen dieser Arbeit wurde PSL erstmals im Kontext einer Konzeptgraph-basierten Wissensgraph-Konstruktion eingesetzt.
Es wurde ein Set von PSL-Regeln~\tref{sec:text2kg:psl:inference} vorgestellt, welches den Einsatz in diesem Kontext exemplarisch aufzeigt.

Neben der Erweiterbarkeit war auch die Parallelisierbarkeit eines der in~\ref{sec:intro:goals} beschriebenen Ziele.
Durch die Wahl von PSL für die Grapherweiterung wurde diese Anforderung prinzipiell erfüllt.
Das für die Inferenz verwendete ADMM-Verfahren lässt sich gut verteilt lösen.
In~\cite[Kapitel 10]{Boyd2011} wird beschrieben, wie sich ADMM mit verteilten Programmiermodellen, wie z.~B. Map-Reduce oder Pregel, implementieren lässt.
Aufbauend darauf, wird in~\cite{Lubell-Doughtie2013} eine konkrete Map-Reduce-basierte Implementation vorgestellt.
PSL ist prinzipiell also für den Einsatz in Cluster-Umgebungen geeignet.
In der im Rahmen dieser Arbeit erstellten Implementation wurde jedoch die Single-Machine ADMM-Implementation der PSL-Referenzimplementation verwendet, da diese für einen Prototypen völlig ausreichend ist.

\section{Ausblick}%
\label{sec:conclusion:todo}
